\chapter{Sets and Structures}\label{ch:b2:structures}

This opening chapter sharpens the foundations laid in the Year 1
volume into working tools of the trade: the calculus of sets and
quotients, the comparison of infinite sets (\omterm{def:b2:structures:countable}{countability},
Cantor--Bernstein), and the structural theory of groups and rings ---
Lagrange's theorem, the symmetric group with its signature, \omterm{def:b2:structures:ideal}{ideals} and
the Chinese remainder theorem. Everything here is used relentlessly in
the rest of the book: the signature builds the determinant
(\cref{ch:b2:linalg}), \omterm{def:b2:structures:quotientring}{quotient rings} drive arithmetic, and
\omterm{def:b2:structures:countable}{countability} underlies both topology and probability.

\section{Sets, maps, quotients}

We use freely the language of sets, maps, and equivalence and \omterm{def:b2:structures:generated}{order}
relations set up in the Year 1 volume. Two upgrades deserve a proper
statement.

\begin{proposition}[Images and preimages of families]\label{prop:b2:structures:images}
Let $f \colon E \to F$ and let $(A_i)_{i \in I}$, $(B_j)_{j \in J}$
be families of subsets of $E$, resp.\ $F$. Then
\[
f^{-1}\Bigl(\bigcup_j B_j\Bigr) = \bigcup_j f^{-1}(B_j),
\qquad
f^{-1}\Bigl(\bigcap_j B_j\Bigr) = \bigcap_j f^{-1}(B_j),
\qquad
f^{-1}(F \setminus B) = E \setminus f^{-1}(B),
\]
\[
f\Bigl(\bigcup_i A_i\Bigr) = \bigcup_i f(A_i),
\qquad
f\Bigl(\bigcap_i A_i\Bigr) \subseteq \bigcap_i f(A_i)
\quad (\text{equality for injective } f).
\]
\end{proposition}

\begin{proof}
Each identity is an unwinding of definitions; for instance $x \in
f^{-1}(\bigcap B_j) \iff f(x) \in B_j$ for all $j$ $\iff x \in
f^{-1}(B_j)$ for all $j$. The image identities and the failure of
equality in the intersection case (with the injectivity fix) were
proved in the Year 1 volume for two sets; the arguments are identical
for families.
\end{proof}

\begin{example}[Where the image inclusion is strict]\label{ex:b2:structures:imagestrict}
Take $f \colon \R \to \R$, $f(x) = x^2$, with $A_1 =
\intcc{-1}{0}$ and $A_2 = \intcc{0}{1}$. Then
\[
f(A_1 \cap A_2) = f(\{0\}) = \{0\},
\qquad
f(A_1) \cap f(A_2) = \intcc{0}{1} \cap \intcc{0}{1} =
\intcc{0}{1} :
\]
the inclusion of \cref{prop:b2:structures:images} is as strict
as can be --- the two preimage points $\pm x$ of a common value
live in different $A_i$. Injectivity is exactly what forbids
this splitting, which is why preimages (which never merge
points) satisfy all four identities unconditionally while
images lose the one about intersections. Rule of thumb for the
whole book: push \emph{preimages} through set operations
freely; handle images with care.
\end{example}

\begin{definition}[Quotient set]\label{def:b2:structures:quotient}
Let $\mathcal{R}$ be an equivalence relation on $E$. The
\emph{quotient set}\index{quotient set} $E/\mathcal{R}$ is the set of
equivalence classes; the surjection $\pi \colon E \to E/\mathcal{R}$,
$x \mapsto \mathrm{cl}(x)$, is the \emph{canonical
projection}\index{canonical projection}.

\emph{Universal property (factorization):} if $f \colon E \to F$ is
\emph{compatible} with $\mathcal{R}$ (i.e.\ $x \mathbin{\mathcal{R}}
y \implies f(x) = f(y)$), there is exactly one map $\overline f
\colon E/\mathcal{R} \to F$ with $f = \overline f \circ \pi$.
\end{definition}

\begin{proof}[Proof of the universal property]
Uniqueness: the requirement $f = \overline f \circ \pi$ reads
\[
\overline f\bigl(\mathrm{cl}(x)\bigr) = f(x)
\qquad (x \in E),
\]
and since $\pi$ is surjective, every element of
$E/\mathcal{R}$ is some $\mathrm{cl}(x)$: the values of
$\overline f$ are all forced. Existence: take the display as
the \emph{definition} of $\overline f$; it is unambiguous
precisely by compatibility --- if $\mathrm{cl}(x) =
\mathrm{cl}(y)$, then $x \mathbin{\mathcal{R}} y$, so $f(x) =
f(y)$ and the two candidate values agree --- and it factorizes
$f$ by construction. Note the division of labour: surjectivity
of $\pi$ gives uniqueness, compatibility gives existence.
\end{proof}

\begin{example}
$\Z/n\Z$ is the quotient of $\Z$ by congruence modulo $n$; the
well-definedness checks of the Year 1 volume were instances of the
universal property. Quotients turn ``compatible constructions on
representatives'' into honest maps --- we use this constantly below.
\end{example}

\section{Countability and cardinality}

\begin{definition}[Equipotence, countability]\label{def:b2:structures:countable}
Two sets are \emph{equipotent}\index{equipotence} when a bijection
joins them. A set is \emph{countable}\index{countable set} when it is
equipotent to $\N$ (some authors include finite sets; we say
\emph{at most countable} for ``finite or countable'').
\end{definition}

\begin{proposition}[Stability properties]\label{prop:b2:structures:countablestable}
\begin{enumerate}
  \item Every infinite subset of $\N$ is \omterm{def:b2:structures:countable}{countable}; a set is at most
        \omterm{def:b2:structures:countable}{countable} iff it injects into $\N$ iff it is empty or a
        surjective image of $\N$.
  \item $\N \times \N$ is \omterm{def:b2:structures:countable}{countable}; a product of two at most
        \omterm{def:b2:structures:countable}{countable sets} is at most \omterm{def:b2:structures:countable}{countable}.
  \item An at most \omterm{def:b2:structures:countable}{countable} union of at most \omterm{def:b2:structures:countable}{countable sets} is at
        most \omterm{def:b2:structures:countable}{countable}.
  \item $\Z$ and $\Q$ are \omterm{def:b2:structures:countable}{countable}.
\end{enumerate}
\end{proposition}

\begin{proof}
(1) List an infinite $A \subseteq \N$ by repeated minima: $a_0 =
\min A$, $a_{k+1} = \min\,(A \setminus \{a_0, \dots, a_k\})$
(nonempty since $A$ is infinite); the map $k \mapsto a_k$ is strictly
increasing, injective, and surjective onto $A$ (every $a \in A$
exceeds only finitely many elements of $A$, so it is reached). If $E$
injects into $\N$ via $\varphi$, then $E$ is \omterm{def:b2:structures:countable}{equipotent} to
$\varphi(E) \subseteq \N$: finite or \omterm{def:b2:structures:countable}{countable}. If $s \colon \N \to
E$ is surjective, then $x \mapsto \min s^{-1}(\{x\})$ injects $E$
into $\N$.

(2) The map $(p, q) \mapsto 2^p(2q + 1) - 1$ is a bijection $\N^2
\to \N$ (every positive integer has a unique odd--even split
$2^p m$ with $m$ odd, by unique factorization). Products: compose
injections.

(3) Given sets $E_n$ with surjections $s_n \colon \N \to E_n$
(harmless when some $E_n$ is finite: repeat values), the map $(n, k)
\mapsto s_n(k)$ is a surjection from the \omterm{def:b2:structures:countable}{countable} $\N^2$ onto
$\bigcup E_n$.

(4) $\Z = \N \cup (-\N^*)$: \omterm{def:b2:structures:countable}{countable} union. $\Q$ is a surjective
image of $\Z \times \N^*$ (the fraction map), hence at most
\omterm{def:b2:structures:countable}{countable}, and infinite.
\end{proof}

\begin{example}[A pairing function, computed]\label{ex:b2:structures:pairing}
The bijection $(p, q) \mapsto 2^p(2q + 1) - 1$ of the proof
deserves to be seen at work. Its first values:
\[
\begin{array}{c|ccccc}
 & q = 0 & q = 1 & q = 2 & q = 3 & q = 4\\
\hline
p = 0 & 0 & 2 & 4 & 6 & 8\\
p = 1 & 1 & 5 & 9 & 13 & 17\\
p = 2 & 3 & 11 & 19 & 27 & 35\\
p = 3 & 7 & 23 & 39 & 55 & 71
\end{array}
\]
Row $p$ collects the integers $n$ for which $n + 1$ is exactly
divisible by $2^p$: every natural number appears exactly once.
Decoding is as explicit as encoding: for $n = 43$, factor $n + 1
= 44 = 2^2\cdot 11 = 2^2(2\cdot5 + 1)$, so $(p, q) = (2, 5)$.
The closing insight: \omterm{def:b2:structures:countable}{countability} proofs are often
\emph{algorithms} in disguise --- here, ``factor out the twos''.
\end{example}

\begin{example}[The algebraic numbers are countable]\label{ex:b2:structures:algebraic}
A complex number is \emph{algebraic} when it annihilates some
nonzero polynomial with rational coefficients. The set
$\overline\Q$ of algebraic numbers is \omterm{def:b2:structures:countable}{countable}: polynomials of
degree $\leq d$ over $\Q$ inject into $\Q^{d+1}$, a finite
product of \omterm{def:b2:structures:countable}{countable sets}
(\cref{prop:b2:structures:countablestable}~(2)); the union over
$d$ enumerates the nonzero rational polynomials as $P_0, P_1,
P_2, \dots$; each $P_k$ has finitely many roots; and
\[
\overline\Q = \bigcup_{k \in \N}\ \{\text{roots of } P_k\}
\]
is a \omterm{def:b2:structures:countable}{countable} union of finite sets
(\cref{prop:b2:structures:countablestable}~(3)), infinite since
it contains $\Q$. Combined with the uncountability of $\R$
(\cref{thm:b2:structures:cantor} below), this proves --- without
exhibiting a single one --- that transcendental numbers exist
and form an uncountable majority: Cantor's counting argument of
1874, existence by cardinality alone.
\end{example}

\begin{theorem}[Cantor; uncountability of $\R$]\label{thm:b2:structures:cantor}
\begin{enumerate}
  \item For every set $E$, there is no surjection $E \to
        \mathcal{P}(E)$.
  \item $\R$ is \emph{not} \omterm{def:b2:structures:countable}{countable}.
\end{enumerate}
\end{theorem}

\begin{proof}
(1) was proved in the Year 1 volume (the diagonal set $D = \{x : x
\notin f(x)\}$).

(2) Suppose $(x_n)_{n \in \N}$ enumerates $\R$. Build nested segments
$I_0 \supseteq I_1 \supseteq \dots$ with $\abs{I_n} = 3^{-n}$ and
$x_n \notin I_n$: split the current segment into three closed thirds;
at least one third avoids $x_n$ (a point meets at most two of the
three). The nested-segments theorem (adjacent endpoints) provides
$\ell \in \bigcap_n I_n$; but $\ell = x_N$ for some $N$, and $x_N
\notin I_N$: contradiction.
\end{proof}

\begin{theorem}[Cantor--Bernstein]\label{thm:b2:structures:cantorbernstein}
If $E$ injects into $F$ and $F$ injects into $E$, then $E$ and $F$
are \omterm{def:b2:structures:countable}{equipotent}.\index{Cantor--Bernstein theorem}
\end{theorem}

\begin{proof}
Let $f \colon E \to F$ and $g \colon F \to E$ be injections. For
each point (of $E$ or $F$), trace its \emph{ancestor chain} of
successive preimages, $x \mapsto g^{-1}(x) \mapsto
f^{-1}(g^{-1}(x)) \mapsto \dots$ --- each step is defined as long as
the current point lies in the image of the relevant injection, and
is then unique by injectivity. Three mutually exclusive fates:
the chain stops at a point of $E \setminus g(F)$ (\emph{origin in
$E$}), stops at a point of $F \setminus f(E)$ (\emph{origin in
$F$}), or never stops. This partitions $E = E_E \cup E_F \cup
E_\infty$ and $F = F_E \cup F_F \cup F_\infty$ according to the
origin.

Now observe: $f$ maps $E_E$ \emph{onto} $F_E$ --- the chain of
$f(x)$ is the chain of $x$ prefixed with one step, so origins match;
and every $y \in F_E$ has a chain with at least one step (its origin
lies in $E$), so $y = f(x)$ with $x \in E_E$. The same argument
gives bijections $f \colon E_\infty \to F_\infty$ and $g \colon F_F
\to E_F$. Gluing,
\[
h(x) =
\begin{cases}
f(x) & \text{if } x \in E_E \cup E_\infty,\\
g^{-1}(x) & \text{if } x \in E_F,
\end{cases}
\]
is a bijection from $E$ onto $F = F_E \cup F_\infty \cup F_F$: it is
bijective piecewise, and the three target pieces are disjoint.
\end{proof}

\begin{example}\label{ex:b2:structures:cbexample}
$\intoo{0}{1}$ and $\intcc{0}{1}$ are \omterm{def:b2:structures:countable}{equipotent}: the identity
injects one way, $x \mapsto \frac{x + 1}{3}$ the other; the theorem
manufactures the (necessarily discontinuous) bijection. Likewise
$\R$, $\intoo{0}{1}$ (via $\tanh$-type bijections) and
$\mathcal{P}(\N)$ (binary expansions, \cref{exo:b2:structures:3})
are all \omterm{def:b2:structures:countable}{equipotent}: ``the cardinality of the continuum''.
\end{example}

\begin{example}[The segment and the square]\label{ex:b2:structures:squareseg}
$\intcc{0}{1}$ and $\intcc{0}{1}^2$ are \omterm{def:b2:structures:countable}{equipotent} --- dimension
is invisible to cardinality. One injection is trivial: $x
\mapsto (x, 0)$. For the other, send $(x, y)$ to the real whose
decimal digits interleave those of $x$ and $y$,
\[
(0.x_1x_2x_3\dots,\ 0.y_1y_2y_3\dots)
\;\longmapsto\; 0.x_1y_1x_2y_2x_3y_3\dots,
\]
choosing for each coordinate the expansion that does not end in
all $9$'s: with that convention the digits of the image
determine those of $x$ and $y$, so the map is injective (it need
not be surjective --- images never have, say, odd-position
digits eventually $9$ --- and that is fine).
Cantor--Bernstein (\cref{thm:b2:structures:cantorbernstein})
assembles a genuine bijection. Continuity, of course, is
hopeless: a continuous bijection between them is impossible ---
the metric chapters explain why (connectedness distinguishes
the line from the plane, \cref{ch:b2:metric}).
\end{example}

\section{Groups}

\begin{definition}[Generated subgroup; order]\label{def:b2:structures:generated}
Let $G$ be a group and $A \subseteq G$. The subgroup \emph{generated}
by $A$, written $\langle A \rangle$, is the smallest subgroup
containing $A$ --- concretely, all finite products of elements of $A$
and their inverses. A group is \emph{cyclic}\index{cyclic group} when
generated by one element: $\langle a\rangle = \{a^k : k \in \Z\}$.
The \emph{order}\index{order!of an element} of $a \in G$ is
$\operatorname{ord}(a) = \abs{\langle a \rangle}$ (possibly
infinite); when finite, it is the least $n \geq 1$ with $a^n = e$,
and $a^k = e \iff \operatorname{ord}(a) \mid k$.
\end{definition}

\begin{proof}[Proof of the order characterization]
If some $a^m = e$ with $m \geq 1$, let $n \geq 1$ be least with $a^n
= e$. The elements $e, a, \dots, a^{n-1}$ are pairwise distinct
($a^{i} = a^{j}$ with $0 \leq i < j < n$ gives $a^{j-i} = e$,
contradicting minimality), and every $a^k$ reduces to one of them by
Euclidean division $k = nq + r$: $\langle a\rangle$ has exactly $n$
elements, and $a^k = a^r = e \iff r = 0 \iff n \mid k$. If no power
is trivial, all $a^k$ ($k \in \Z$) are distinct (same division
argument) and the \omterm{def:b2:structures:generated}{order} is infinite.
\end{proof}

\begin{theorem}[Lagrange]\label{thm:b2:structures:lagrange}
Let $G$ be a finite group and $H$ a subgroup. Then $\abs H$ divides
$\abs G$.\index{Lagrange's theorem} In particular the \omterm{def:b2:structures:generated}{order} of every
element divides $\abs G$, and $a^{\abs G} = e$ for all $a \in G$.
\end{theorem}

\begin{proof}
The relation $x \sim y \iff x^{-1}y \in H$ is an equivalence
(reflexive: $e \in H$; symmetric: inverses; transitive: products).
The class of $x$ is the \emph{left coset} $xH = \{xh : h \in H\}$,
and $h \mapsto xh$ is a bijection $H \to xH$ (inverse $y \mapsto
x^{-1}y$): all classes have $\abs H$ elements. Classes partition $G$
(the general partition theorem of the Year 1 volume), so $\abs G =
\abs H \times (\text{number of cosets})$. For an element: apply this
to $H = \langle a\rangle$; then $a^{\abs G} = (a^{\operatorname{ord}
a})^{\abs G / \operatorname{ord} a} = e$.
\end{proof}

\begin{example}[Cosets in action: $A_3$ inside $\mathfrak{S}_3$]\label{ex:b2:structures:cosets}
Take $G = \mathfrak{S}_3$ (\omterm{def:b2:structures:generated}{order} $6$) and $H = A_3 =
\{\mathrm{id},\ (1\,2\,3),\ (1\,3\,2)\}$. The left cosets are
\[
H = \{\mathrm{id},\ (1\,2\,3),\ (1\,3\,2)\},
\qquad
(1\,2)H = \{(1\,2),\ (2\,3),\ (1\,3)\} :
\]
two classes of three elements partitioning $G$, exactly as the
count $\abs G = \abs H \times (\text{number of cosets})$
demands --- and visibly the partition into even and odd
permutations. Note $(1\,3)H = (1\,2)H$ although $(1\,3) \neq
(1\,2)$: cosets are \emph{classes}, not labelled by their
representatives, and $x^{-1}y \in H$ is the only legitimate
comparison. This two-class picture is the general one for the
signature: $A_n$ and its lone companion coset split
$\mathfrak{S}_n$ in half, which is how the weekend problem
counts reachable puzzle positions.
\end{example}

\begin{example}\label{ex:b2:structures:lagrangeapp}
Two immediate dividends. \emph{Groups of prime \omterm{def:b2:structures:generated}{order} are \omterm{def:b2:structures:generated}{cyclic}:}
if $\abs G = p$ is prime and $a \neq e$, then
$\operatorname{ord}(a)$ divides $p$ and is not $1$, so it is $p$:
$\langle a\rangle = G$. \emph{The subgroup lattice of $\Z/12\Z$:}
by \cref{prop:b2:structures:cyclic} below, there is exactly one
subgroup per divisor of $12$ --- \omterm{def:b2:structures:generated}{orders} $1, 2, 3, 4, 6, 12$,
\omterm{def:b2:structures:generated}{generated} respectively by $\overline 0$, $\overline 6$, $\overline
4$, $\overline 3$, $\overline 2$, $\overline 1$. The closing
caution: the \emph{converse} of Lagrange fails in general --- $A_4$
has \omterm{def:b2:structures:generated}{order} $12$ but no subgroup of \omterm{def:b2:structures:generated}{order} $6$, as we prove in this
chapter's weekend problem (\cref{pb:b2:structures:1}, question 14).
Lagrange restricts the possible \omterm{def:b2:structures:generated}{orders}; it does not promise them.
\end{example}

\begin{omfigure}
\begin{tikzpicture}[scale=1.05, every node/.style={font=\small}]
  \node (G) at (0,3)  {$\Z/12\Z = \langle\overline1\rangle$};
  \node (A) at (-1.6,2) {$\langle\overline2\rangle$};
  \node (B) at (1.6,2)  {$\langle\overline3\rangle$};
  \node (C) at (-1.6,1) {$\langle\overline4\rangle$};
  \node (D) at (1.6,1)  {$\langle\overline6\rangle$};
  \node (E) at (0,0)  {$\{\overline0\}$};
  \draw[omDef] (G) -- (A); \draw[omDef] (G) -- (B);
  \draw[omDef] (A) -- (C); \draw[omDef] (A) -- (D);
  \draw[omDef] (B) -- (D); \draw[omDef] (C) -- (E);
  \draw[omDef] (D) -- (E);
  \node[omThm, right=2pt] at (0.1,3) {\ order $12$};
  \node[omThm, left=2pt]  at (-2.1,2) {order $6$\ };
  \node[omThm, right=2pt] at (2.1,2) {\ order $4$};
  \node[omThm, left=2pt]  at (-2.1,1) {order $3$\ };
  \node[omThm, right=2pt] at (2.1,1) {\ order $2$};
\end{tikzpicture}

{\small The subgroup lattice of $\Z/12\Z$: one subgroup per
divisor of $12$ (\cref{prop:b2:structures:cyclic}), with an edge
when one contains the other with prime index. Inclusions run
\emph{against} divisibility of the generator: $\langle\overline
4\rangle \subseteq \langle\overline2\rangle$ because $4$ is a
multiple of $2$.}
\end{omfigure}

\begin{proposition}[Cyclic groups]\label{prop:b2:structures:cyclic}
Let $G = \langle a \rangle$ be \omterm{def:b2:structures:generated}{cyclic} of \omterm{def:b2:structures:generated}{order} $n$.
\begin{enumerate}
  \item $G$ is isomorphic to $(\Z/n\Z, +)$, via $\overline k \mapsto
        a^k$.
  \item Every subgroup of $G$ is \omterm{def:b2:structures:generated}{cyclic}; for each divisor $d \mid n$
        there is exactly one subgroup of \omterm{def:b2:structures:generated}{order} $d$, namely $\langle
        a^{n/d}\rangle$.
  \item $a^k$ generates $G$ if and only if $\gcd(k, n) = 1$: $G$ has
        $\varphi(n)$ generators (Euler's function\index{Euler's
        totient function}).
\end{enumerate}
\end{proposition}

\begin{proof}
(1) The map $k \mapsto a^k$ from $\Z$ onto $G$ is compatible with
congruence mod $n$ ($a^{k} = a^{k'} \iff n \mid k - k'$, by the \omterm{def:b2:structures:generated}{order}
characterization); the universal property
(\cref{def:b2:structures:quotient}) yields a well-defined bijective
morphism from $\Z/n\Z$.

(2) Let $H \leq G$ be nontrivial and $m$ least $\geq 1$ with $a^m \in
H$. Euclidean division shows $H = \langle a^m\rangle$ (for $a^k \in
H$: $k = mq + r$ forces $a^r \in H$, so $r = 0$), and $m \mid n$
(divide $n$ by $m$: $a^{n \bmod m} \in H$). Then $\abs H = n/m$;
taking $m = n/d$ realizes each divisor $d$. Uniqueness: any subgroup
of \omterm{def:b2:structures:generated}{order} $d$ is, by the above, of the form $\langle a^m \rangle$
with $n/m = d$ --- so $m = n/d$ is forced and the subgroup is
determined.

(3) We claim $\operatorname{ord}(a^k) = \frac{n}{\gcd(k, n)}$.
Write $d = \gcd(k, n)$. For any $m \geq 1$, the \omterm{def:b2:structures:generated}{order}
characterization of \cref{def:b2:structures:generated} gives the
chain of equivalences
\[
(a^k)^m = e
\iff n \mid km
\iff \frac{n}{d} \,\Big|\, \frac{k}{d}\,m
\iff \frac{n}{d} \,\Big|\, m ,
\]
the last step by Gauss's lemma, since $\frac nd$ and $\frac kd$
are coprime. The least such $m$ is $\frac nd$:
$\operatorname{ord}(a^k) = \frac n{\gcd(k,n)}$, which equals $n$
iff $\gcd(k, n) = 1$. There are $\varphi(n)$ such classes $k$
modulo $n$.
\end{proof}

\section{The symmetric group}

\begin{definition}\label{def:b2:structures:sn}
$\mathfrak{S}_n$ is the group of permutations of $\intint{1}{n}$
(\omterm{def:b2:structures:generated}{order} $n!$). A \emph{cycle}\index{cycle} $(a_1\,a_2\,\cdots\,a_k)$
maps $a_1 \mapsto a_2 \mapsto \dots \mapsto a_k \mapsto a_1$ and
fixes everything else; $k$ is its \emph{length}, a $2$-cycle is a
\emph{transposition}\index{transposition}. Two cycles are
\emph{disjoint} when their
supports (non-fixed points) are.
\end{definition}

\begin{theorem}[Cycle decomposition]\label{thm:b2:structures:cycles}
Every permutation $\sigma \neq \mathrm{id}$ is a product of pairwise
disjoint \omterm{def:b2:structures:sn}{cycles}, uniquely up to the order of the factors. Disjoint
\omterm{def:b2:structures:sn}{cycles} commute, and $\operatorname{ord}(\sigma)$ is the lcm of the
lengths.
\end{theorem}

\begin{proof}
Consider the ``orbit'' relation on the support of $\sigma$: $x \sim y$
iff $y = \sigma^k(x)$ for some $k \in \Z$ --- an equivalence
relation. Each class $\{x, \sigma(x), \dots, \sigma^{k-1}(x)\}$
(finite, so the iterates \omterm{def:b2:structures:sn}{cycle} back --- the first repetition must
return to $x$ by injectivity) carries the \omterm{def:b2:structures:sn}{cycle} $(x\ \sigma(x)\
\cdots\ \sigma^{k-1}(x))$, and $\sigma$ is the product of these
\omterm{def:b2:structures:sn}{cycles}: on each orbit, only the corresponding \omterm{def:b2:structures:sn}{cycle} acts.
Uniqueness: any \omterm{def:b2:structures:sn}{disjoint-cycle} factorization reproduces exactly the
orbits (the \omterm{def:b2:structures:sn}{cycle} through $x$ must be $(x\ \sigma(x)\ \cdots)$).
Disjoint \omterm{def:b2:structures:sn}{cycles} commute since they move disjoint points; the \omterm{def:b2:structures:generated}{order}
statement follows because $\sigma^m = \mathrm{id}$ iff each \omterm{def:b2:structures:sn}{cycle}'s
$m$-th power is (disjointness), iff each length divides $m$.
\end{proof}

\begin{example}[Cycle type as a census]\label{ex:b2:structures:cycletype}
How many permutations of $\mathfrak{S}_9$ have the \omterm{def:b2:structures:sn}{cycle} type
$(4, 3, 2)$ --- one $4$-cycle, one $3$-cycle, one \omterm{def:b2:structures:sn}{transposition}?
Choose the supports and the cyclic orders:
\[
\frac{9!}{4\cdot 3\cdot 2}
= \frac{362\,880}{24} = 15\,120 :
\]
list the nine symbols in a row ($9!$ ways), bracket the first
four, next three, last two into \omterm{def:b2:structures:sn}{cycles}, and divide by the
rotations inside each bracket ($4$, $3$ and $2$ of them) which
give the same permutation. (Distinct \omterm{def:b2:structures:sn}{cycle} \emph{lengths} here,
so no further division; equal lengths would also require
dividing by the permutations of the equal brackets.) Every such
permutation has \omterm{def:b2:structures:generated}{order} $\operatorname{lcm}(4,3,2) = 12$ and
signature $(-1)^3(-1)^2(-1)^1 = +1$
(\cref{thm:b2:structures:cycles} and the signature theorem
below). One partition of $9$, one conjugacy class, one census
--- the combinatorics of $\mathfrak{S}_n$ is the arithmetic of
partitions.
\end{example}

\begin{theorem}[Signature]\label{thm:b2:structures:signature}
There is exactly one group morphism $\varepsilon \colon
\mathfrak{S}_n \to \{\pm 1\}$ (for $n \geq 2$) taking the value $-1$
on \omterm{def:b2:structures:sn}{transpositions}: the \emph{signature}\index{signature}. Moreover
$\varepsilon(\sigma) = (-1)^{I(\sigma)}$ where $I(\sigma)$ is the
number of \emph{inversions} (pairs $i < j$ with $\sigma(i) >
\sigma(j)$), a $k$-cycle has signature $(-1)^{k-1}$, and the
\emph{alternating group} $A_n = \ker\varepsilon$ has \omterm{def:b2:structures:generated}{order}
$\frac{n!}{2}$.
\end{theorem}

\begin{proof}
\emph{Existence.} For $\sigma \in \mathfrak{S}_n$ set
\[
\varepsilon(\sigma)
= \prod_{1 \leq i < j \leq n}
\frac{\sigma(j) - \sigma(i)}{j - i} .
\]
The factors' absolute values multiply to $1$ (the unordered pairs
$\{\sigma(i), \sigma(j)\}$ run over all pairs), so
$\varepsilon(\sigma) = (-1)^{I(\sigma)} \in \{\pm1\}$. Morphism: for
$\sigma, \tau$,
\[
\varepsilon(\sigma\tau)
= \prod_{i<j} \frac{\sigma(\tau(j)) - \sigma(\tau(i))}{j - i}
= \prod_{i<j} \frac{\sigma(\tau(j)) - \sigma(\tau(i))}{\tau(j) -
\tau(i)} \cdot \prod_{i<j} \frac{\tau(j) - \tau(i)}{j - i}
= \varepsilon(\sigma)\,\varepsilon(\tau),
\]
the middle product being $\varepsilon(\sigma)$ after reindexing by
the pairs $\{\tau(i), \tau(j)\}$ (each unordered pair appears once,
and numerator and denominator flip sign together). A \omterm{def:b2:structures:sn}{transposition}
$\tau = (a\,b)$ with $a < b$ has an odd number of inversions;
counted exactly: the inverted pairs $(i, j)$, $i < j$, with
$\tau(i) > \tau(j)$ are
\[
(a, j) \ \text{for } a < j < b, \qquad
(i, b) \ \text{for } a < i < b, \qquad
(a, b) \ \text{itself},
\]
that is $(b - a - 1) + (b - a - 1) + 1 = 2(b - a) - 1$ of them,
odd. (Alternatively: check $(1\,2)$ directly, with one
inversion, and conjugate --- conjugates have equal signature
since $\varepsilon$ is a morphism to an abelian group.) Hence
$\varepsilon((a\,b)) = (-1)^{2(b-a)-1} = -1$.

\emph{Uniqueness.} \omterm{def:b2:structures:sn}{Transpositions} generate $\mathfrak{S}_n$ (any
\omterm{def:b2:structures:sn}{cycle} $(a_1\cdots a_k) = (a_1\,a_k)(a_1\,a_{k-1})\cdots(a_1\,a_2)$,
and \cref{thm:b2:structures:cycles} finishes); a morphism to
$\{\pm1\}$ is determined by its values on generators.

\emph{Consequences.} The \omterm{def:b2:structures:sn}{cycle} identity above writes a $k$-cycle as
$k - 1$ \omterm{def:b2:structures:sn}{transpositions}: signature $(-1)^{k-1}$. $A_n$: the morphism
$\varepsilon$ is surjective (\omterm{def:b2:structures:sn}{transpositions} exist for $n \geq 2$),
and the two ``cosets'' $A_n$ and $(1\,2)A_n$ are \omterm{def:b2:structures:countable}{equipotent} and
partition $\mathfrak{S}_n$ (Lagrange's argument): $\abs{A_n} =
\frac{n!}{2}$.
\end{proof}

\begin{example}\label{ex:b2:structures:snexample}
$\sigma = \begin{pmatrix} 1&2&3&4&5&6\\ 3&6&5&4&1&2 \end{pmatrix}
= (1\,3\,5)(2\,6)$: \omterm{def:b2:structures:generated}{order} $\operatorname{lcm}(3,2) = 6$, signature
$(-1)^{2}\cdot(-1)^{1} = -1$. The signature is the fastest
parity check on shuffles --- and the engine of the determinant in
\cref{ch:b2:linalg}.
\end{example}

\begin{example}[Three roads to one sign]\label{ex:b2:structures:threeroads}
Let $\sigma \in \mathfrak{S}_5$ send $1, 2, 3, 4, 5$ to $3, 5, 4,
1, 2$. \emph{Via \omterm{def:b2:structures:sn}{cycles}:} $1 \mapsto 3 \mapsto 4 \mapsto 1$ and $2
\mapsto 5 \mapsto 2$, so $\sigma = (1\,3\,4)(2\,5)$ and
$\varepsilon(\sigma) = (-1)^{2}(-1)^{1} = -1$. \emph{Via
inversions:} in the value list $3, 5, 4, 1, 2$ the out-of-order
pairs are $(3,1)$, $(3,2)$, $(5,4)$, $(5,1)$, $(5,2)$, $(4,1)$,
$(4,2)$: seven of them, and $(-1)^7 = -1$. \emph{Via
\omterm{def:b2:structures:sn}{transpositions}:} $\sigma = (1\,4)(1\,3)(2\,5)$, three factors,
$(-1)^3 = -1$. Three computations, one parity: the uniqueness in
\cref{thm:b2:structures:signature} guarantees that no bookkeeping
scheme can ever make them disagree --- which is exactly what makes
$\varepsilon$ usable as an invariant (see the weekend problem).
\end{example}

\begin{remark}[Where the signature goes from here]
The signature is the seed of three later harvests: it builds the
determinant and its product rule in \cref{ch:b2:linalg}; it powers
parity invariants for combinatorial puzzles (this chapter's weekend
problem solves the fifteen puzzle with it); and the alternating
groups $A_n$ it defines become central in the Year 3 volume, where
their simplicity for $n \geq 5$ explains why degree-$5$ equations
have no solution in radicals.
\end{remark}

\section{Rings, ideals, quotients}

\begin{definition}[Ideal]\label{def:b2:structures:ideal}
Let $A$ be a commutative ring. An \emph{ideal}\index{ideal} $I
\subseteq A$ is an additive subgroup such that $a x \in I$ for all
$a \in A$, $x \in I$. Kernels of ring morphisms are ideals; $I = A$
iff $1 \in I$ iff $I$ contains a unit. The ideal \emph{\omterm{def:b2:structures:generated}{generated}} by
$x$ is $xA = \{xa\}$ (a \emph{principal} ideal).
\end{definition}

\begin{theorem}[{Ideals of $\Z$ and of $K[X]$}]\label{thm:b2:structures:principal}
Every \omterm{def:b2:structures:ideal}{ideal} of $\Z$ is $n\Z$ for a unique $n \in \N$; every \omterm{def:b2:structures:ideal}{ideal} of
$K[X]$ ($K$ a field) is $P\,K[X]$ for a unique monic (or zero) $P$.
Consequently gcd's exist in both rings with Bézout relations: $x\Z +
y\Z = \gcd(x,y)\Z$, and likewise for polynomials.
\end{theorem}

\begin{proof}
For $\Z$ this was the subgroup theorem of the Year 1 volume (an \omterm{def:b2:structures:ideal}{ideal}
is in particular a subgroup, and $n\Z$ is an \omterm{def:b2:structures:ideal}{ideal}). For $K[X]$: let
$I \neq \{0\}$ be an \omterm{def:b2:structures:ideal}{ideal} and $P \in I$ nonzero of minimal degree,
normalized monic. For $F \in I$, Euclidean division $F = PQ + R$
gives $R = F - PQ \in I$ with $\deg R < \deg P$: minimality forces $R
= 0$, so $I = P\,K[X]$. Uniqueness: two monic generators divide each
other. The Bézout statements are the equality of the \omterm{def:b2:structures:ideal}{ideal} $x\Z +
y\Z$ (resp.\ its polynomial analogue) with the principal \omterm{def:b2:structures:ideal}{ideal} of the
gcd --- the very definition of gcd used in Year 1, now recognized as
a statement about \omterm{def:b2:structures:ideal}{ideals}.
\end{proof}

\begin{example}[A polynomial gcd, two ways]\label{ex:b2:structures:polygcd}
Compute $\gcd(X^3 - 1,\ X^2 - 1)$ in $\Q[X]$. \emph{By Euclid:}
\[
X^3 - 1 = X\,(X^2 - 1) + (X - 1),
\qquad
X^2 - 1 = (X + 1)(X - 1) + 0 ,
\]
so the gcd is $X - 1$, and back-substitution gives the Bézout
relation
\[
X - 1 = 1\cdot(X^3 - 1) - X\cdot(X^2 - 1).
\]
\emph{By \omterm{def:b2:structures:ideal}{ideals}:} the \omterm{def:b2:structures:ideal}{ideal} $(X^3 - 1)\Q[X] + (X^2 - 1)\Q[X]$ is
principal (\cref{thm:b2:structures:principal}); it contains $X -
1$ (the display) and is contained in $(X - 1)\Q[X]$ (both
generators vanish at $1$, hence are multiples of $X - 1$): the
monic generator is $X - 1$. The closing insight: the \omterm{def:b2:structures:ideal}{ideal}
viewpoint identifies the gcd \emph{without dividing} --- common
roots locate the \omterm{def:b2:structures:ideal}{ideal}, and Euclid merely certifies it.
\end{example}

\begin{definition}[Quotient ring $\Z/n\Z$, revisited]\label{def:b2:structures:quotientring}
For an \omterm{def:b2:structures:ideal}{ideal} $I$ of $A$, the relation $x \sim y \iff x - y \in I$ is
an equivalence compatible with $+$ and $\times$; the \omterm{def:b2:structures:quotient}{quotient set}
$A/I$ inherits a ring structure --- the \emph{quotient
ring}\index{quotient ring} --- making $\pi \colon A \to A/I$ a
morphism with kernel $I$. For $A = \Z$, $I = n\Z$ this is the
$\Z/n\Z$ of the Year 1 volume, now with its universal property: any
morphism killing $I$ factors through $A/I$.
\end{definition}

\begin{theorem}[Chinese remainder theorem, ring form]\label{thm:b2:structures:crt}
If $\gcd(m, n) = 1$, the map
\[
\Z/mn\Z \longrightarrow \Z/m\Z \times \Z/n\Z,
\qquad
\overline{x} \longmapsto (x \bmod m,\; x \bmod n)
\]
is a ring isomorphism.\index{Chinese remainder theorem} Consequently
$\varphi(mn) = \varphi(m)\varphi(n)$ for coprime $m, n$, and
\[
\varphi(n) = n \prod_{p \mid n} \Bigl(1 - \frac 1p\Bigr)
\quad (p \text{ prime}).
\]
\end{theorem}

\begin{proof}
The map is a well-defined ring morphism (compatibilities are
immediate). Injectivity: $x \equiv 0$ mod $m$ and mod $n$ with
$\gcd(m,n) = 1$ forces $mn \mid x$ (Gauss). Surjectivity: both sides
have $mn$ elements, so injectivity suffices (finite equal
cardinalities) --- or explicitly: from a B\'ezout relation $um +
vn = 1$, the class of
\[
x = b\,um + a\,vn
\]
maps to $(a \bmod m,\ b \bmod n)$, since $vn = 1 - um \equiv 1
\pmod m$ makes $x \equiv a \pmod m$, and symmetrically mod $n$
--- the recipe used numerically in
\cref{ex:b2:structures:crtinverse}. Units correspond to pairs of units (a product ring's
units are the pairs of units), so $\varphi(mn) =
\varphi(m)\varphi(n)$. For a prime power, $\varphi(p^k) = p^k -
p^{k-1}$ (the non-units mod $p^k$ are the multiples of $p$);
multiplicativity assembles the product formula.
\end{proof}

\begin{example}[Inverting the Chinese isomorphism]\label{ex:b2:structures:crtinverse}
Take $m = 8$, $n = 9$. The inverse of the isomorphism is made
explicit by the two \emph{idempotents}: seek $u \equiv 1 \pmod 8$,
$u \equiv 0 \pmod 9$ and $v \equiv 0 \pmod 8$, $v \equiv 1 \pmod
9$. From $u = 9k \equiv 1 \pmod 8$: $k \equiv 1$, so $u = 9$; from
$v = 8k \equiv 1 \pmod 9$: $-k \equiv 1$, $k \equiv 8$, so $v =
64$. Then the class of $x = 9a + 64b$ modulo $72$ is the unique
solution of $x \equiv a \pmod 8$, $x \equiv b \pmod 9$: for $a =
3$, $b = 5$ one gets $27 + 320 = 347 \equiv 59 \pmod{72}$ ---
exactly the intermediate value found by substitution in
\cref{exo:b2:structures:8}. The closing insight: $u$ and $v$
satisfy $u + v \equiv 1$, $uv \equiv 0$, $u^2 \equiv u$, $v^2
\equiv v$ modulo $72$; they are the images of $(1, 0)$ and $(0,
1)$, and every Chinese decomposition is at bottom a decomposition
of $1$ into orthogonal idempotents.
\end{example}

\begin{theorem}[Euler; Fermat revisited]\label{thm:b2:structures:euler}
The units of $\Z/n\Z$ form a group of \omterm{def:b2:structures:generated}{order} $\varphi(n)$; hence for
$\gcd(a, n) = 1$:
\[
a^{\varphi(n)} \equiv 1 \pmod n
\qquad (\text{Euler's theorem\index{Euler's theorem}}),
\]
and Fermat's little theorem is the case $n = p$ prime, now one line
from Lagrange.
\end{theorem}

\begin{proof}
The invertible classes are exactly those of integers coprime to $n$
(Year 1 volume): $\varphi(n)$ of them, forming a group under
multiplication. Lagrange (\cref{thm:b2:structures:lagrange}):
every element to the power of the group \omterm{def:b2:structures:generated}{order} is the identity.
\end{proof}

\begin{example}[A unit group without a generator]\label{ex:b2:structures:units15}
The group $(\Z/15\Z)^*$ has $\varphi(15) = \varphi(3)\varphi(5)
= 8$ elements. Is it \omterm{def:b2:structures:generated}{cyclic}? Compute \omterm{def:b2:structures:generated}{orders} using the Chinese
isomorphism $(\Z/15\Z)^* \simeq (\Z/3\Z)^* \times (\Z/5\Z)^*$
(a unit mod $15$ is a pair of units): the factors have \omterm{def:b2:structures:generated}{orders}
$2$ and $4$, so every element's \omterm{def:b2:structures:generated}{order} divides
$\operatorname{lcm}(2, 4) = 4 < 8$ --- no element generates.
Concretely:
\[
2^4 = 16 \equiv 1, \qquad
4^2 = 16 \equiv 1, \qquad
7^4 \equiv 1, \qquad
11^2 = 121 \equiv 1, \qquad
14^2 \equiv 1 \pmod{15} :
\]
\omterm{def:b2:structures:generated}{orders} $4, 2, 4, 2, 2$ and never $8$. Contrast with
\cref{exo:b2:structures:10}: $(\Z/p\Z)^*$ \emph{is} \omterm{def:b2:structures:generated}{cyclic} for
$p$ prime, because there the unit group sits inside a field.
Euler's theorem still applies with exponent $\varphi(15) = 8$,
but the true universal exponent here is $4$ --- Euler is an
upper bound, not always the sharp one.
\end{example}

\begin{definition}[Algebra]\label{def:b2:structures:algebra}
A \emph{$K$-algebra}\index{algebra} is a $K$-vector space $A$ with a
ring structure whose multiplication is $K$-bilinear. Examples:
$K[X]$, $\mathcal{M}_n(K)$, $\mathcal{L}(E)$, function spaces
$\mathcal{F}(X, K)$, $\C$ as an $\R$-algebra. Morphisms of algebras
are linear ring morphisms; the \emph{evaluation} $P \mapsto P(u)$
from $K[X]$ to $\mathcal{L}(E)$ (or $\mathcal{M}_n(K)$) is the
central example, driving \cref{ch:b2:reduction}.
\end{definition}

\begin{example}[An evaluation morphism and its kernel]\label{ex:b2:structures:evaluation}
Take $A = \begin{pmatrix}0 & 1\\ 0 & 0\end{pmatrix}$ and the
evaluation $\varepsilon_A \colon \R[X] \to \mathcal{M}_2(\R)$,
$P \mapsto P(A)$. Since $A^2 = 0$,
\[
P(A) = P(0)\,I + P'(0)\,A =
\begin{pmatrix} P(0) & P'(0)\\ 0 & P(0)\end{pmatrix},
\]
(only the constant and linear terms of $P$ survive). Hence
$\ker\varepsilon_A = \{P : P(0) = P'(0) = 0\} = X^2\,\R[X]$: a
principal \omterm{def:b2:structures:ideal}{ideal}, exactly as \cref{thm:b2:structures:principal}
predicts, \omterm{def:b2:structures:generated}{generated} by the monic $X^2$ of least degree in the
kernel --- the \emph{minimal polynomial} of $A$, star of
\cref{ch:b2:reduction}. The image is the two-dimensional
commutative \omterm{def:b2:structures:algebra}{algebra} $\{aI + bA\}$: evaluation morphisms shrink
the infinite-dimensional $\R[X]$ onto small, computable \omterm{def:b2:structures:algebra}{algebras}.
\end{example}

\begin{remark}[Perspectives: three melodies to listen for]\label{rem:b2:structures:perspectives}
Three structural ideas from this chapter recur throughout the
volume, each time in heavier orchestration. \emph{Factorization
through a quotient} (\cref{def:b2:structures:quotient}): it
builds $\Z/n\Z$ here, defines maps on solution sets of linear
systems in \cref{ch:b2:linalg}, and silently underlies every
``well defined on classes'' argument. \emph{Invariants}: the
signature is a morphism to $\{\pm1\}$ that no legal move can
dodge --- the same logic gives the determinant's product rule
(\cref{ch:b2:linalg}), the trace's similarity invariance, and
the conserved quantities of \cref{ch:b2:diffeq}.
\emph{Counting against a structure}: Lagrange counts through
cosets, dimension counts through bases
(\cref{ch:b2:linalg}), multiplicity counts through polynomial
degrees (\cref{ch:b2:reduction}); whenever a bound looks
miraculous, some partition or grading is doing the counting.
\end{remark}

\begin{remark}[Common pitfalls]\label{rem:b2:structures:pitfalls}
Four classics. (i) A map on a quotient must be checked
\emph{well defined}: ``$\overline x \mapsto$ (formula on $x$)''
is legitimate only if the formula is constant on classes ---
the compatibility of \cref{def:b2:structures:quotient}, not a
formality. (ii) $\operatorname{ord}(ab) =
\operatorname{lcm}(\operatorname{ord}a, \operatorname{ord}b)$
is \emph{false} in general, even for commuting elements ($a$
and $a^{-1}$); \cref{exo:b2:structures:4} gives the correct
coprime-and-commuting statement, and disjoint \omterm{def:b2:structures:sn}{cycles} the
correct permutation version. (iii) \omterm{def:b2:structures:countable}{Countability} survives
\omterm{def:b2:structures:countable}{countable} \emph{unions} and finite \emph{products}, but not
\omterm{def:b2:structures:countable}{countable} products: $\{0,1\}^{\N}$ is uncountable
(\cref{exo:b2:structures:3}) although each factor has two
elements. (iv) Cantor--Bernstein needs only injections both
ways, but the bijection it builds is usually discontinuous and
non-explicit --- do not expect a formula
(\cref{ex:b2:structures:cbexample}).
\end{remark}

\begin{remark}[Where this chapter is used]
Almost everywhere. The signature builds determinants
(\cref{ch:b2:linalg}); the evaluation morphism $P \mapsto P(u)$
and the principal \omterm{def:b2:structures:ideal}{ideals} of $K[X]$ produce minimal polynomials and
the kernel decompositions of \cref{ch:b2:reduction}; \omterm{def:b2:structures:countable}{countability}
is the stage on which \cref{ch:b2:proba} performs (probability on
\omterm{def:b2:structures:countable}{countable} spaces) and the reason topology keeps producing \omterm{def:b2:structures:countable}{countable}
dense sets (\cref{ch:b2:metric}). The quotient construction $A/I$
is redeployed in the Year 3 volume to build fields $K[X]/(P)$ and,
from them, Galois theory: the universal property proved here is
used there word for word.
\end{remark}

\section{Exercises}

\begin{exercise}[$\star$]\label{exo:b2:structures:1}
Which of the following sets are \omterm{def:b2:structures:countable}{countable}? The set of finite subsets
of $\N$; the set of \emph{all} subsets of $\N$; $\R \setminus \Q$;
the set of polynomials with rational coefficients; the set of
sequences of $0$'s and $1$'s that are eventually zero.
\end{exercise}

\begin{exercise}[$\star$]\label{exo:b2:structures:2}
In $\mathfrak{S}_7$, let $\sigma = (1\,4\,2\,6)(3\,5)$ and $\tau =
(2\,3\,7)$. Compute $\sigma\tau$ and $\tau\sigma$ in \omterm{def:b2:structures:sn}{disjoint-cycle}
form, the \omterm{def:b2:structures:generated}{orders} and signatures of all four permutations, and
$\sigma^{2026}$.
\end{exercise}

\begin{exercise}[$\star$]\label{exo:b2:structures:3}
Construct explicit injections showing that $\mathcal{P}(\N)$,
$\intcc{0}{1}$ and the set $\{0,1\}^{\N}$ of binary sequences are
pairwise \omterm{def:b2:structures:countable}{equipotent} \emph{(binary expansions both ways;
Cantor--Bernstein absorbs the double-representation nuisance)}.
\end{exercise}

\begin{exercise}[$\star$]\label{exo:b2:structures:4}
Let $G$ be a group and $a, b \in G$ commuting elements of finite
coprime \omterm{def:b2:structures:generated}{orders} $m$ and $n$. Prove that $\operatorname{ord}(ab) =
mn$. Show by an example in $\mathfrak{S}_3$ that commutation is
essential.
\end{exercise}

\begin{exercise}[$\star\star$]\label{exo:b2:structures:5}
Let $G$ be a finite group of even \omterm{def:b2:structures:generated}{order}. Prove that $G$ contains an
element of \omterm{def:b2:structures:generated}{order} $2$. \emph{(Pair each element with its inverse;
count the self-paired ones.)}
\end{exercise}

\begin{exercise}[$\star\star$]\label{exo:b2:structures:6}
Prove that $A_n$ ($n \geq 3$) is \omterm{def:b2:structures:generated}{generated} by the $3$-cycles.
\emph{(A product of two \omterm{def:b2:structures:sn}{transpositions} is a $3$-cycle or a product
of two $3$-cycles.)}
\end{exercise}

\begin{exercise}[$\star\star$]\label{exo:b2:structures:7}
Determine all group morphisms: from $(\Q, +)$ to $(\Z, +)$; from
$(\Z/n\Z, +)$ to $(\Z/m\Z, +)$ \emph{(count them:
$\gcd(m,n)$)}; from $(\Q, +)$ to $(\Q_+^*, \times)$.
\end{exercise}

\begin{exercise}[$\star\star$]\label{exo:b2:structures:8}
Using the Chinese remainder theorem, compute $\varphi(360)$, find
all $x$ with $x \equiv 3 \pmod 8$, $x \equiv 5 \pmod 9$ and $x
\equiv 2 \pmod 5$, and compute the last two digits of $3^{2026}$
\emph{(Euler mod $100$; beware: work mod $4$ and mod $25$)}.
\end{exercise}

\begin{exercise}[$\star\star\star$]\label{exo:b2:structures:9}
Prove that a finite integral domain is a field. Deduce that
$\Z/n\Z$ is a field iff $n$ is prime (again).
\end{exercise}

\begin{exercise}[$\star\star\star$]\label{exo:b2:structures:10}
(A classic) Let $K$ be a field and $G$ a \emph{finite} subgroup of
$(K^*, \times)$. Prove that $G$ is \omterm{def:b2:structures:generated}{cyclic}.
\emph{Hint: let $m$ be the maximal \omterm{def:b2:structures:generated}{order} among elements of $G$;
show every element's \omterm{def:b2:structures:generated}{order} divides $m$ (using
\cref{exo:b2:structures:4} on suitable coprime parts), so all of $G$
satisfies $x^m = 1$; count roots of $X^m - 1$.}
In particular $(\Z/p\Z)^*$ is \omterm{def:b2:structures:generated}{cyclic}.
\end{exercise}

\begin{exercise}[$\star\star\star$]\label{exo:b2:structures:11}
Prove that the group $(\Q, +)$ is not \omterm{def:b2:structures:generated}{cyclic}, and worse: it is not
even finitely \omterm{def:b2:structures:generated}{generated}. Prove on the other hand that every finitely
\omterm{def:b2:structures:generated}{generated} subgroup of $(\Q, +)$ is \omterm{def:b2:structures:generated}{cyclic}.
\end{exercise}

\begin{exercise}[$\star\star$]\label{exo:b2:structures:12}
(Dedekind's criterion) Prove that every infinite set contains a
\omterm{def:b2:structures:countable}{countable} subset, and deduce that a set $E$ is infinite if and only
if it is \omterm{def:b2:structures:countable}{equipotent} to a proper subset of itself. \emph{(For the
direct implication, shift a \omterm{def:b2:structures:countable}{countable} subset by one step; for the
converse, recall the pigeonhole principle.)}
\end{exercise}

\section{Problem: The Fifteen Puzzle}

The fifteen puzzle is a $4 \times 4$ tray holding fifteen sliding
tiles numbered $1$ to $15$ and one empty cell; a move slides one of
the tiles adjacent to the empty cell into it. In the 1890s Sam Loyd
popularized the puzzle by offering \$1000 to anyone who could
exchange the tiles $14$ and $15$ and return every other tile to its
place. Nobody ever collected, and this weekend problem proves both
halves of the reason: the signature of
\cref{thm:b2:structures:signature} forbids Loyd's exchange, and
--- the harder, constructive half --- \emph{everything} the
signature allows is genuinely solvable. The full statement is the
Johnson--Story theorem (1879).

\begin{omfigure}
\begin{tikzpicture}[scale=0.72]
  \draw[omDef, thick] (0,0) grid (4,4);
  \draw[omDef, very thick] (0,0) rectangle (4,4);
  \foreach \k in {1,...,15} {
    \pgfmathtruncatemacro\row{int((\k-1)/4)}
    \pgfmathtruncatemacro\col{int(mod(\k-1,4))}
    \node at (\col+0.5, 3.5-\row) {$\k$};
  }
  \node[below, font=\small] at (2,-0.2) {solved};
\end{tikzpicture}
\qquad
\begin{tikzpicture}[scale=0.72]
  \draw[omDef, thick] (0,0) grid (4,4);
  \draw[omDef, very thick] (0,0) rectangle (4,4);
  \foreach \k in {1,...,13} {
    \pgfmathtruncatemacro\row{int((\k-1)/4)}
    \pgfmathtruncatemacro\col{int(mod(\k-1,4))}
    \node at (\col+0.5, 3.5-\row) {$\k$};
  }
  \node[omThm] at (1.5,0.5) {$15$};
  \node[omThm] at (2.5,0.5) {$14$};
  \node[below, font=\small] at (2,-0.2) {Loyd's bounty};
\end{tikzpicture}

{\small The solved configuration and Sam Loyd's $14$--$15$
configuration. The \$1000 question: can legal slides turn the right
tray into the left one?}
\end{omfigure}

\begin{problem}[{Weekend problem --- the Johnson--Story
solvability theorem}]
\label{pb:b2:structures:1}
Number the cells $1$ to $16$ in reading order (left to right, top
to bottom), so that cell $k$ sits in row $i$ and column $j$ with $k
= 4(i - 1) + j$. Cell $16$ (bottom right) is the \emph{home} of
the empty cell; we treat the empty cell as a sixteenth tile,
written $b$ and identified with the number $16$. A
\emph{configuration} is a bijection $\sigma \colon \intint1{16}
\to \intint1{16}$, cell $\mapsto$ content; the \emph{solved}
configuration is $\sigma = \mathrm{id}$. Throughout,
$\varepsilon$ is the signature of
\cref{thm:b2:structures:signature} and two cells are
\emph{adjacent} when they share an edge of the tray.

\medskip
\noindent\textbf{Part I --- Configurations, moves, signatures.}
\begin{enumerate}
  \item Justify that the configurations are exactly the elements
        of $\mathfrak{S}_{16}$, so there are $16! =
        20\,922\,789\,888\,000$ of them, and that the number of
        legal moves from a given configuration is $2$, $3$ or $4$,
        according to whether the empty cell lies in a corner, on
        an edge, or in the interior.
  \item Let $\sigma$ be a configuration, $p = \sigma^{-1}(16)$ the
        cell of the blank, and $c$ a cell adjacent to $p$. Show
        that sliding the tile of $c$ into $p$ produces the
        configuration $\sigma' = \sigma \circ \tau$ with $\tau =
        (p\ c)$, and deduce that every move flips the signature:
        $\varepsilon(\sigma') = -\varepsilon(\sigma)$.
  \item Checkerboard the tray: $\chi(k) = (-1)^{i+j}$ for the cell
        $k$ in row $i$, column $j$. Show that every move flips
        $\chi(\text{cell of the blank})$, and deduce that a
        sequence of moves returning the blank to its starting cell
        has even length.
  \item Show that
        \[
        I(\sigma) = \varepsilon(\sigma)\,
        \chi\bigl(\sigma^{-1}(16)\bigr)
        \]
        is invariant under every legal move, and compute
        $I(\mathrm{id})$.
\end{enumerate}

\medskip
\noindent\textbf{Part II --- Loyd's bounty: the invariant at
work.}
\begin{enumerate}[resume]
  \item Loyd's configuration $\sigma_L$ agrees with the solved one
        except that cells $14$ and $15$ hold tiles $15$ and $14$.
        Compute $I(\sigma_L)$ and conclude that no sequence of
        moves links $\sigma_L$ to the solved configuration:
        Loyd's \$1000 was never in danger.
  \item Show that exactly half of all configurations satisfy $I =
        +1$: $\abs{\{\sigma : I(\sigma) = +1\}} = 16!/2$.
        \emph{(For a fixed blank cell, pair configurations by
        composing with one fixed \omterm{def:b2:structures:sn}{transposition} of two other
        cells.)}
  \item Show that every move is undone by a legal move, that
        ``$\sigma'$ is reachable from $\sigma$ by legal moves'' is
        an equivalence relation, and that the class $R$ of the
        solved configuration satisfies $R \subseteq \{I = +1\}$.
        Conclude that there are at least two classes.
  \item Suppose the blank is home: $\sigma(16) = 16$. Show that
        $I(\sigma) = \varepsilon(\rho)$ where $\rho \in
        \mathfrak{S}_{15}$ is the restriction of $\sigma$ to the
        cells $1, \dots, 15$, and that any configuration can be
        carried by legal moves to one with the blank home.
        Conclude: to prove $R = \{I = +1\}$ it suffices to realize
        every \emph{even} permutation of the fifteen non-home
        cells by a sequence of moves starting and ending with the
        blank home.
\end{enumerate}

\medskip
\noindent\textbf{Part III --- Blank tours and the program
group.} A \emph{program} is a finite sequence of legal moves,
started from a configuration with the blank home, whose final
configuration again has the blank home. Its \emph{effect} is the
permutation $\pi$ of the cells defined by: the content of cell $x$
ends in cell $\pi(x)$.
\begin{enumerate}[resume]
  \item Show that a program run from $\sigma$ ends at $\sigma
        \circ \pi^{-1}$; that running two programs in succession
        composes their effects; and that the set $H$ of all
        effects is a subgroup of $\mathfrak{S}_{15}$ (permutations
        of the cells $1, \dots, 15$) contained in the alternating
        group $A_{15}$.
  \item (The elementary tour) From the blank at home, slide the
        blank around the bottom-right $2 \times 2$ block: cells
        $16 \to 12 \to 11 \to 15 \to 16$. Show the effect is the
        $3$-cycle $(11\ 12\ 15)$, and that the reverse tour gives
        $(11\ 15\ 12)$. Both lie in $H$.
  \item (The grand tour) Verify that
        \[
        16 \to 15 \to 14 \to 13 \to 9 \to 5 \to 1 \to 2 \to 3
        \to 4 \to 8 \to 7 \to 6 \to 10 \to 11 \to 12 \to 16
        \]
        is a closed walk through all sixteen cells (adjacent steps
        only), and that its effect is the $15$-cycle
        \[
        \zeta = (15\ 12\ 11\ 10\ 6\ 7\ 8\ 4\ 3\ 2\ 1\ 5\ 9\ 13\
        14) .
        \]
        Writing $x_0 = 15$, $x_1 = 12$, $x_2 = 11$, \dots, $x_{14}
        = 14$ for its \omterm{def:b2:structures:sn}{cycle} \omterm{def:b2:structures:generated}{order}, check that the reverse
        elementary tour of question 10 is exactly $(x_0\ x_1\
        x_2)$.
  \item Prove the conjugation formula in any $\mathfrak{S}_n$: for
        a permutation $g$ and a $3$-cycle,
        \[
        g\,(a\ b\ c)\,g^{-1} = \bigl(g(a)\ g(b)\ g(c)\bigr),
        \]
        and note that $H$, being a group, is closed under
        conjugation by its own elements.
  \item Deduce that $H$ contains all fifteen \emph{consecutive}
        $3$-cycles of the grand tour:
        \[
        s_t = (x_t\ x_{t+1}\ x_{t+2}) \qquad (t \in \Z/15\Z,
        \text{ indices mod } 15).
        \]
\end{enumerate}

\medskip
\noindent\textbf{Part IV --- Generating the alternating group.}
\begin{enumerate}[resume]
  \item (Lemma A) Let $s$ and $t$ be $3$-cycles whose supports
        share exactly two points, say supports $\{a, b, c\}$ and
        $\{b, c, d\}$. Show that, after replacing $s$ or $t$ by
        its inverse if necessary (which changes nothing to the
        \omterm{def:b2:structures:generated}{generated} subgroup), the product $st$ is a double
        \omterm{def:b2:structures:sn}{transposition}; show that $A_4$ contains no subgroup of
        \omterm{def:b2:structures:generated}{order} $6$ \emph{(a subgroup of index $2$ contains every
        square; count the $3$-cycles among squares)}; and conclude
        that $\langle s, t\rangle$ is the whole alternating group
        of the four letters $\{a, b, c, d\}$.
  \item (Lemma B) Let $X$ be a set of $k \geq 4$ letters, $w
        \notin X$, and let $G$ be a subgroup of some
        $\mathfrak{S}_n$ containing every even permutation of $X$
        and one $3$-cycle $(u\ v\ w)$ with $u, v \in X$. Show that
        for all distinct $a, b \in X$ there is an \emph{even}
        permutation $g$ of $X$ with $g(u) = a$, $g(v) = b$, and
        deduce $(a\ b\ w) \in G$.
  \item Deduce that the group $G$ of Lemma B contains every even
        permutation of $X \cup \{w\}$ \emph{(use
        \cref{exo:b2:structures:6}: the $3$-cycles generate)}.
        Then, chaining Lemmas A and B along the consecutive
        $3$-cycles $s_0, s_1, \dots, s_{12}$ of question 13,
        prove that $\langle s_0, \dots, s_{12}\rangle = A_{15}$.
  \item Conclude that $H = A_{15}$: \emph{every even rearrangement
        of the fifteen tiles is achievable by a program}, and $H$
        has $15!/2 = 653\,837\,184\,000$ elements.
  \item (The Johnson--Story theorem, 1879) Assemble questions 6,
        7, 8 and 17: the configurations reachable from the solved
        one are \emph{exactly} the $16!/2 =
        10\,461\,394\,944\,000$ configurations with $I = +1$;
        and reachability has exactly \emph{two} classes, the class
        of the solved configuration and the class of Loyd's
        $\sigma_L$. \emph{(For the second point, relabel the tiles
        $14$ and $15$: show $\sigma \mapsto (14\ 15) \circ
        \sigma$ maps move sequences to move sequences and
        exchanges $\{I = +1\}$ with $\{I = -1\}$.)}
\end{enumerate}

\medskip
\noindent\textbf{Part V --- Criteria, variants, and the view
from above.}
\begin{enumerate}[resume]
  \item (The practical criterion) Read the fifteen tiles in
        reading order of their cells, skipping the blank, and let
        $N$ be the number of inversions of this list; let $r$ be
        the row of the blank counted from the \emph{bottom}. Show
        that $I(\sigma) = (-1)^{N + r + 1}$, so that $\sigma$ is
        solvable if and only if $N + r$ is odd.
  \item (Group actions) An \emph{action} of a group $G$ on a set
        $X$ is a map $G \times X \to X$, $(g, x) \mapsto g \cdot
        x$, with $e \cdot x = x$ and $g \cdot (h \cdot x) =
        (gh) \cdot x$; the \emph{orbit} of $x$ is $G \cdot x$, and
        the action is \emph{free} when $g \cdot x = x$ forces $g =
        e$. Show that $h \cdot \sigma = \sigma \circ h^{-1}$
        defines a free action of $H$ on the set of blank-home
        configurations, that its orbits are exactly the classes of
        mutual reachability by programs, and recover from the
        orbit count that these configurations split into exactly
        $15!\,/\,\abs H = 2$ classes.
  \item (The $3 \times 3$ obstruction) Show that the $3 \times 3$
        board admits \emph{no} closed walk visiting every cell
        exactly once: the grand-tour strategy of Part III fails
        for the eight puzzle. \emph{(Checkerboard the nine
        cells.)}
  \item (The repair) On the $3 \times 3$ board with cells $1$ to
        $9$ in reading order and home $9$: compute the effects of
        the perimeter tour $9 \to 8 \to 7 \to 4 \to 1 \to 2 \to 3
        \to 6 \to 9$ (a $7$-cycle $\zeta'$ fixing the center $5$)
        and of the corner tour $9 \to 6 \to 5 \to 8 \to 9$ (a
        $3$-cycle through the center). Conjugating the latter by
        the powers of $\zeta'$ and chaining Lemmas A and B, prove
        that the eight puzzle's program group is all of $A_8$,
        hence that exactly $9!/2 = 181\,440$ of the $9! =
        362\,880$ configurations are solvable.
  \item (A poor board) Now let the board be a single \omterm{def:b2:structures:sn}{cycle} of $n
        \geq 4$ cells carrying $n - 1$ tiles. Show that the cyclic
        order of the tiles is invariant, that each reachability
        class has exactly $n(n - 1)$ configurations \emph{(the
        classes are the orbits of a \omterm{def:b2:structures:generated}{cyclic group} of \omterm{def:b2:structures:generated}{order}
        $\operatorname{lcm}(n, n-1) = n(n-1)$)}, and that there
        are $(n - 2)!$ classes --- for $n \geq 5$ far more than
        $2$: on a thin board the parity invariant captures almost
        nothing, and geometry rules.
  \item Two verdicts by the criterion of question 19: the fully
        reversed tray (tiles $15, 14, \dots, 1$ in cells $1$
        to $15$, blank home) and the tray with the blank in cell
        $1$ followed by the tiles $15, 14, \dots, 1$ in cells $2$
        to $16$. Which one is solvable?
  \item (Synthesis) The proof has two independent pillars: an
        \emph{invariant} ($I$, built from the signature morphism)
        showing at most half the configurations are reachable, and
        an \emph{explicit generation} theorem ($H = A_{15}$)
        showing at least half are. In one sentence each, say where
        the following entered: the morphism property of
        $\varepsilon$; Lagrange's theorem; the generation of $A_n$
        by $3$-cycles; conjugation. State the meta-principle in
        one line.
\end{enumerate}
\end{problem}
